\documentclass[12pt]{article}

\usepackage[a4paper,total={6in, 10in}]{geometry}
\usepackage{xcolor}

\begin{document}

\title{Research Proposal}
\author{Jonathan Reicher \\ Shachar Itzhaki}
\maketitle

\newcommand{\todo}[1]{\textcolor{red}{TODO: #1}}
\setlength{\parskip}{0.1cm}

\section*{Problem Statement}
In the intersection between the worlds of mathematical theory, formal logic,
and programming, exists a species of tools called proof assistants.
These are programming languages, meant not for building software, but for
building mathematical proofs.
Proof assistants, contrary to pen and paper, check formalized proofs for
correctness and soundness.
These systems have existed for many years, but recently, one in particular has
been gaining traction.

Lean (\todo{cite?}) is an interactive proof assistant geared for a general math
audience.
Two things that make it especially appealing are its large ecosystem and its
intuitive syntax.
Despite that, there are some challenging aspects of adopting Lean, and one in
particular is how complex its semantics and type system are, and how.

We propose building a system that will reduce the friction of using Lean by
being more ``forgiving'' and less formal. The system will allow users to write
down proofs in way that humans can read and understand as written, with parts of
the proof being translated to Lean, to help a user start formalizing their
proof in Lean.
This workflow is in accordance to how researchers tend to use Lean, as a formal
proof in Lean often comes as a supplementary part of a paper proof inside a
research paper.

\section*{Existing Solutions}
Mention open ai’s paper that adi sent
Mention aristotel

\section*{Preliminary Work}
Started from Yuval’s domino example. Tried to encode the domain and the solution
in LEAN.
embed proof.md and a matching lean code that we expect

\section*{References}

\end{document}
